\chapter{Methodology and Design}
\label{chapter:methodology_and_design}

\section{Proposed Analytical Framework}

To bridge the gap between traditional flood modeling and the operational demands of a Reactive Digital Twin, this research proposes a highly integrated, multi-layered framework. Designed to ingest and process the high-velocity data streams characteristic of flash floods—such as the 2024 Valencia DANA—the architecture is built upon three core pillars:

\begin{enumerate}
    \item \textbf{Discrete Space Formalization ($m:n\text{-CA}^k$)}: Leveraging cellular automata to represent complex, non-linear hydraulic interactions via discrete transition rules, offering a computationally efficient alternative to traditional Shallow Water Equation (SWE) solvers.
    \item \textbf{Hexagonal Architecture (Technological Decoupling)}: Separating the core simulation physics from external dependencies (GIS parsing, MQTT messaging, logging) to ensure modularity and seamless adaptability across different hardware and data infrastructures.
    \item \textbf{Distributed Telemetry}: Implementing an MQTT-based protocol layer to transform the standalone execution engine into a distributed node, enabling real-time, low-latency synchronization with remote visualization pipelines and emergency decision-support systems.
\end{enumerate}

This framework shifts flood simulation away from traditional offline batch-processing toward a continuous, concurrent loop of data ingestion, parallelized execution, and telemetry broadcast, establishing the essential operational infrastructure for a Reactive Digital Twin.


\subsection{The Multilayer Conceptualization ($m:n-CA^k$)}

The cornerstone of the simulation logic is the adoption of the Multilayer Cellular Automata ($m:n-CA^k$) formalism. Traditional 
Cellular Automata (CA) typically operate on a single state per cell, which is insufficient for representing the heterogeneous variables 
of a river basin. The $m:n-CA^k$ generalization allows for a multidimensional representation where the environment is discretized into 
a stack of $n$ computational layers.

\paragraph{Mathematical Formalization}

Formally, the simulation environment is defined as a specialized cellular space where each cell $c$ at a given coordinates $(i, j)$ is 
characterized by a state vector $S_{i,j}$. The system is represented by the tuple:

$$A = \langle R, E, V, N, \Psi \rangle$$

Where:

\begin{itemize}
    \item $R$ is the $m$-dimensional cell space.
    \item $E$ is the set of $n$ layers that define the global state and the physical environment of the system.
    \item $V$ represents the vicinity function.
    \item $N$ represents the nucleus function.
    \item $\Psi$ is the combination or evolution function that governs the fluid dynamics.
\end{itemize}


\subsubsection{Cell Space ($R$) and Layer Definition ($E$)}

The simulation space $R$ is defined on a regular two-dimensional grid $R \subset \mathbb{Z}^2$, where each cell is 
identified by its spatial coordinates $(i, j)$. The physical size of each cell is determined by the spatial resolution 
$\Delta x$.

The system consists of a set $E = \{E_1, E_2, E_3, E_4, E_5, E_6\}$ of $n=6$ computational layers. One of the fundamental characteristics 
of the $m:n-CA^k$ model is that the layers may differ in their geometric and temporal dimensions.

\begin{enumerate}
    \item $E_1$: Topobathymetry ($Z_{topo-bathy}$). A static 2D layer of dimension $X \times Y$. It combines the Digital 
    Elevation Model (DEM) with the bathymetry of the watercourses, providing the absolute elevation of the riverbed: 
    $E_1 \subset \mathbb{R}^{X \times Y}$.
    \item $E_2$: Surface Roughness ($\mu_{soil}$). A static 2D layer of dimension $X \times Y$. Although derived from discrete 
    Land Cover (LC) classifications via a State Transfer Function ($LC \xrightarrow{f} \mu_{soil}$), this continuous layer of 
    Manning's $n$ coefficients is the one that strictly persists during execution to govern horizontal friction: 
    $E_2 \subset \mathbb{R}^{X \times Y}$.
    \item $E_3$: Water Depth ($WD$). A 2D dynamic layer of dimension $X \times Y$. It represents the primary physical state of the cell 
    at time $t$, storing the surface water depth in meters: 
    $E_3 \subset \mathbb{R}^{X \times Y}$.
    \item $E_4$: Rainfall ($RF$). A 3D dynamic layer of dimension $X \times Y \times T$, explicitly incorporating the temporal 
    domain to inject discrete precipitation volumes ($\Delta V_{in}$) into the cell column during each time step $\Delta t$: 
    $E_4 \subset \mathbb{R}^{X \times Y \times T}$.
    \item $E_5$: Water Movement State ($WMS$). A 2D dynamic logical layer of dimension $X \times Y$ storing discrete states ($\text{int8}$). 
    It tracks whether a cell is hydraulically active (Dynamic vs. Static). This layer acts as a structural optimization shield, allowing 
    the execution engine to skip sub-threshold cells via dynamic bounding boxes: 
    $E_5 \subset \mathbb{N}^{X \times Y}$.
    \item $E_6$: Flood Risk Level ($FR$). A 2D dynamic categorical output layer of dimension $X \times Y$. It classifies the hazard level 
    of each cell at time $t$ by evaluating a heuristic threshold function over the current water depth ($WD \xrightarrow{g} FR$), mapping 
    physical outputs into discrete risk levels for real-time hazard mapping: 
    $E_6 \subset \mathbb{N}^{X \times Y}$.
\end{enumerate}


\subsubsection{Vicinity Function ($V$) and Nucleus Function ($N$)}

The \textbf{Vicinity Function} ($V$) determines the subset of adjacent cells that influence the calculation of the future state of a 
given cell. A Moore neighborhood with radius $r=1$ is adopted, defining a set of 8 adjacent cells connected orthogonally and 
diagonally.

$$V(i,j) = \{ (i+u, j+v) \in R \mid \max(|u|, |v|) \le 1 \}$$

The \textbf{Nucleus Function} ($N$) is projected onto the central cell itself to update the dynamic states of the system ($E_3$, $E_5$, 
and $E_6$), mapping the spatio-temporal modifications through local evolution rules: $N(i,j) = \{ (i,j) \}$.


\subsubsection{Intra-layer Transitions: Horizontal Propagation Dynamics}

Intra-layer transitions define the kinematic behavior of surface water flow across the discretized spatial grid.

\paragraph{A. Neighborhood Formalization and Hydraulic Gradients}

To minimize grid-orientation bias, the simulation framework utilizes a standard Moore neighborhood ($k = 8$). For any given cell $c_{i,j}$, the execution engine computes a vector of localized, non-negative directional hydraulic gradients ($S_k$) based on the total hydraulic head—the sum of topobathymetric elevation ($Z_{\text{topo-bathy}}$) and water depth ($WD$)—and the Euclidean distance $L_k$:

$$S_{k} = \max\left(0, \frac{(Z_{\text{topo-bathy}, i,j} + WD_{i,j}) - (Z_{\text{topo-bathy}, k} + WD_{k})}{L_{k}}\right) \cdot \mathbb{I}(WMS_{i,j})$$

where the distance parameter is defined as $L_k = \Delta x$ for cardinal neighbors ($k \in \{1, 2, 3, 4\}$) and $L_k = \Delta x \sqrt{2}$ for diagonal neighbors ($k \in \{5, 6, 7, 8\}$). 

The binary indicator function $\mathbb{I}(WMS_{i,j}) \in \{0, 1\}$ is derived from the Water Movement State layer ($E_5$). If a cell is flagged as hydraulically inactive, its directional gradients are immediately nullified, acting as a computational shield that bypasses unnecessary downstream tensor operations.

\paragraph{B. Multi-Physics Velocity Vector Field Evaluation}

To accommodate variable fluid rheologies, the multi-physics core independently evaluates distinct turbulent and laminar regimes across all eight Moore neighborhood directions, computing a directional velocity component ($v_k$) for each target cell:

\begin{itemize}
    \item \textbf{Turbulent Regime (Manning)}: For friction-dominated clear-water dynamics, the velocity $v_{\text{turb}, k}$ is evaluated using a tensor-optimized variant of Manning's equation, where $\mu_{\text{soil}}$ represents the terrain roughness coefficient from layer $E_2$:
    $$v_{\text{turb}, k} = \frac{1}{\mu_{\text{soil}}} \cdot WD_{i,j}^{2/3} \cdot \sqrt{S_k}$$
    
    \item \textbf{Laminar / Viscous Regime}: For dense sediment or mudflows, the viscous velocity $v_{\text{lam}, k}$ incorporates fluid density ($\rho$), gravity ($g$), dynamic viscosity ($\mu_{\text{fluid}}$), and a tiny regularization parameter ($\epsilon_{\text{reg}} = 10^{-8}$) to prevent division-by-zero errors:
    $$v_{\text{lam}, k} = \frac{\rho \cdot g \cdot WD_{i,j}^2 \cdot S_k}{3\mu_{\text{fluid}} + \epsilon_{\text{reg}}}$$
\end{itemize}

The core resolves the effective directional velocity $v_{\text{approx}, k}$ via a local minimization function, continuously bounding the vector field by the dominant physical constraint:

$$v_{\text{approx}, k} = \min\left(v_{\text{turb}, k}, v_{\text{lam}, k}\right)$$

\paragraph{C. Directional Volumetric Displacement}

To translate linear velocities ($m/s$) into discrete physical mass variations, the engine computes the directional cross-sectional flow 
area $A_{\text{cs}}$ passing through the cell interfaces, defined as the product of the active cell water depth and the grid resolution:

$$A_{\text{cs}} = WD_{i,j} \cdot \Delta x$$

Consequently, the multi-directional Theoretical Outflow Volume vector $\Delta V_{\text{theor}, k}$ over the discrete simulation time 
step $\Delta t$ is formulated as:

$$\Delta V_{\text{theor}, k} = v_{\text{approx}, k} \cdot A_{\text{cs}} \cdot \Delta t$$

\paragraph{D. Non-Linear Mass Conservation and Multiple Flow Direction (MFD) Routing}

This framework utilizes a Multiple Flow Direction (MFD) architecture where routing fractions are governed by the directional velocity field. To guarantee mass conservation and eliminate numerical oscillations, a localized explicit volumetric limiter restricts outflows based on the cell's maximum available volume ($V_{\text{cell}} = WD_{i,j} \cdot \Delta x^2$). 

A dynamic scaling factor ($\alpha$) is computed to bound the cumulative theoretical outflows to the available storage, ensuring strictly conservative final outflows ($\Delta V_{\text{out}, k}$):

$$\alpha = \frac{\min\left(\sum_{k=1}^{8} \Delta V_{\text{theor}, k},\, V_{\text{cell}}\right)}{\sum_{k=1}^{8} \Delta V_{\text{theor}, k}} \quad \implies \quad \Delta V_{\text{out}, k} = \Delta V_{\text{theor}, k} \cdot \alpha$$

The discrete mass balance equation then updates the water depth using these regulated outbound and inbound ($\Delta V_{\text{in}, k}$) volumes:

$$WD_{i,j}^{t+1} = WD_{i,j}^{t} - \left( \frac{\sum_{k=1}^{8} \Delta V_{\text{out}, k}}{\Delta x^2} \right) + \left( \frac{\sum_{k=1}^{8} \Delta V_{\text{in}, k}}{\Delta x^2} \right)$$

Finally, an epsilon floor filter ($\epsilon = 10^{-5}\text{ m}$) is applied to eliminate sub-threshold floating-point precision noise:

$$WD_{i,j}^{t+1} = \begin{cases} 
0.0, & \text{if } WD_{i,j}^{t+1} < \epsilon \\ 
WD_{i,j}^{t+1}, & \text{otherwise} 
\end{cases}$$

\paragraph{E. Dynamic Computational Shielding (Water Movement State)}

At the closure of each temporal iteration, the automaton updates the dynamic states of the Water Movement State layer ($E_5$). The 
activation state transitions between static and dynamic regimes based on a strict flow activity threshold:

$$WMS_{i,j}^{t+1} = \begin{cases} 
1 \, (\text{Dynamic}), & \text{if } \sum_{k=1}^{8} \Delta V_{\text{out}, k} > \epsilon \lor \sum_{k=1}^{8} \Delta V_{\text{in}, k} > \epsilon \\ 
0 \, (\text{Static}), & \text{otherwise} 
\end{cases}$$

Cells that fall below this threshold are flagged as static, allowing the core execution loop to bypass their tensor state operations in 
subsequent iterations, dramatically optimizing computational throughput over dry areas of the river basin.

\paragraph{F. Real-Time Hazard Categorization}

The final phase of the intra-layer loop populates the Flood Risk Level layer ($E_6$). This categorical output layer maps the continuous 
physical dimensions into emergency management metrics via a discrete heuristic transfer function $g$:

$$FR_{i,j}^{t+1} = g\left(WD_{i,j}^{t+1}\right)$$

Where $g$ categorizes the computed water depths into predefined hazard classes (e.g., Low, Medium, High, Severe Risk), generating 
action-ready spatial indicators at every interval $\Delta t$.

\subsubsection{Inter-layer Transitions: Vertical Propagation Dynamics}

Inter-layer transitions govern the vertical, localized interactions between different information layers within a single cell column 
$c_{i,j}$.

\paragraph{A. Atmospheric Flux Integration (Precipitation Injection)}

The primary hydrological driver of the simulation consists of localized, high-intensity meteoric events. Within the $m:n-CA^k$ algebraic 
framework, precipitation is structured as a decoupled vertical flux that injects discrete water volumes directly into the Water Depth 
layer ($E_3$) prior to horizontal routing:

$$WD_{i,j}^{*} = WD_{i,j}^{t} + \left( R_{i,j}^{t} \cdot \Delta t \right)$$

Where $R_{i,j}^{t}$ represents the rainfall intensity ($m/s$) extracted from the spatiotemporal 3D Rainfall tensor ($E_4$) at time $t$.

\paragraph{B. Discrete-to-Continuous Physical Parameterization (Roughness Generation)}

During initialization ($t=0$), categorical land cover data is transformed into a continuous physical boundary by mapping the Land Cover layer onto the persistent Surface Roughness layer ($E_2$) via a tensorial Lookup Table (LUT):

$$\forall c_{i,j} \in R \implies E_{2(i,j)} = \text{LUT}\left(\mathcal{L}_{\text{LandCover}}(i,j)\right)$$

This transformation allows the engine to purge raw land cover data post-initialization. During runtime, the propagation logic directly queries this continuous Manning's layer $E_2$ via parallel tensor gathers, efficiently capturing varying flow resistances across urban districts and transport infrastructure.



\section{Software Architecture: The Hexagonal Approach}
\label{sec:hexagonal-methodology}

A common vulnerability in hydraulic simulators is the tight coupling between physical models and their underlying computational infrastructure (e.g., geospatial file parsers, communication protocols, hardware libraries). To eliminate these dependencies, this framework adopts a Hexagonal Architecture (Ports and Adapters pattern), isolating the pure domain physics of the flood from external data ingestion and transmission mechanisms. This strict decoupling is essential for a Reactive Digital Twin, allowing the system to seamlessly switch from offline HDF5 file parsing to real-time MQTT rainfall streams without modifying a single line of the core hydrodynamic routing code.

\subsection{Core vs. Infrastructure}

The Hexagonal design divides the platform into two distinct regions governed by an inward-pointing dependency rule: dependencies must always target the core domain logic, leaving the computational core completely unaware of external technical implementations.

\paragraph{A. The Core (The "Inside")}
The Core contains the pure scientific and operational logic of the simulation framework:
\begin{itemize}
    \item \textbf{The $m:n\text{-CA}^k$ Engine}: Implements the discrete transition functions (e.g., Manning's equations, MFD routing).
    \item \textbf{The Domain Model}: Defines the multi-layered tensorial representation of the grid, cell states, and hydraulic attributes.
    \item \textbf{Application Services}: Orchestrates the execution of the main simulation loop and temporal synchronization.
\end{itemize}

\paragraph{B. The Infrastructure (The "Outside")}
The Infrastructure encapsulates all technical details and external dependencies required to support the Core:
\begin{itemize}
    \item \textbf{Data Ingestion}: Handles file I/O operations, static storage, and geospatial data parsing (e.g., DEMs).
    \item \textbf{Computation Backends}: Manages hardware-accelerated execution graphs for tensor matrix operations.
    \item \textbf{Telemetry and Messaging}: Implements network protocols (MQTT) and data serialization for real-time visualization.
\end{itemize}
This architectural segregation ensures that technological evolution or shifts in GIS standards within the infrastructure layer do not compromise the scientific validity of the simulation engine.



\section{Distributed Communication Model (MQTT)}
\label{sec:distributed_communication_model}

Real-time flood forecasting demands low-latency, reliable data delivery alongside high mathematical precision. To bridge the computational core with remote visualization and decision-support clients, the framework integrates a distributed telemetry layer governed by the Message Queuing Telemetry Transport (MQTT) protocol. This lightweight, asynchronous publish-subscribe standard minimizes computational overhead and ensures reliable data dissemination over unstable networks during crisis events like the 2024 Valencia DANA.

\subsection{Decoupling Strategy: Publisher-Subscriber Paradigm}
\label{subsec:decoupling_strategy}

To prevent network latency or client-side rendering bottlenecks from throttling the core execution engine, the communication architecture enforces strict decoupling between the engine (Publisher) and remote clients (Subscribers) across three operational dimensions:

\begin{itemize}
    \item \textbf{Spatial and Identity Anonymity}: The simulation core functions as an isolated data publisher, remaining completely agnostic of downstream network topologies or consumer counts. This allows multiple external entities to connect dynamically without adding memory or computational overhead to the main execution thread.
    
    \item \textbf{Temporal Decoupling}: Hydraulic execution remains independent of client consumption rates. Network disturbances or rendering lag are buffered by an intermediate message broker, isolating the simulator and preserving real-time performance guarantees.
    
    \item \textbf{Functional Isolation}: Abstracting network interfaces into outbound data ports ensures that the core physical and mathematical cell state transitions remain unpolluted by transport protocols or session management policies.
\end{itemize}

\subsection{Formal State Distribution Channels and Topic Hierarchy}
\label{subsec:topic_hierarchy}

Following the distributed simulation paradigms established by \cite{FonsecaICasas2013}, state transitions are mapped into structured, event-driven telemetry channels to isolate behavioral model execution from graphical representation. Rooted in a unique scenario identifier, the MQTT topic namespace is segregated into three functional pathways based on transmission lifecycle:

\begin{itemize}
    \item \textbf{Lifecycle and Peer Discovery} (\texttt{FloodSim/\{scenarioName\}/system \\ /handshake}): A low-frequency administrative channel managing operational readiness, peer handshakes, and capability negotiations prior to initialization.
    
    \item \textbf{Structural and Physical State Broadcast} (\texttt{FloodSim/\{scenarioName\} \\ /events}): The primary high-traffic data pipeline handling an initial setup phase (broadcasting domain boundaries and matrix resolutions) followed by a continuous execution phase (broadcasting dynamic water depths and hazard levels).
    
    \item \textbf{Bi-directional Control and Orchestration} (\texttt{FloodSim/\{scenarioName\} \\ /control/events}): An out-of-band feedback loop receiving upstream commands, backpressure metrics, and frame acknowledgments to maintain active flow control between distributed computing nodes.
\end{itemize}

\subsection{Differential State Framing and Temporal Synchronization Protocol}
\label{subsec:temporal_synchronization}

To mitigate bandwidth constraints over distributed networks, the framework replaces full-state data sweeps with a \textit{Differential State Update} approach, transmitting only the cell deltas that drift beyond a predefined physical tolerance. For each simulation step $t$, chronological state coherency across remote visualizers is maintained via a strict four-phase orchestration sequence:

\begin{enumerate}
    \item \textbf{Frame Initialization}: Emits a control boundary signal declaring the onset of a new temporal frame and the total volume of upcoming discrete sub-packets.
    
    \item \textbf{State Modification}: Disperses the physical updates across the network, carrying differential cell coordinates, localized fluid depths, and risk levels.
    
    \item \textbf{Frame Termination}: Transmits a closure signal notifying visualizers that all differential updates for step $t$ have been delivered, safely triggering a scene graph refresh.
    
    \item \textbf{Temporal Synchronization}: Executes a synchronization heartbeat mapping the internal execution clock to the external rendering time base, preventing visual artifacts and temporal drift across nodes.
\end{enumerate}


\section{Data Engineering and Tensor Representation}

The efficiency of a Cellular Automata simulation is fundamentally constrained by its data access patterns. While traditional GIS 
software operates on various file formats and coordinate systems, a high-performance simulator requires a unified, contiguous, and 
multidimensional memory layout. This section details the "Data-as-Interface" philosophy, where geospatial information is abstracted 
into tensors to enable vectorized computation.


\subsection{The Hierarchical Spatiotemporal Framework}
\label{subsec:hierarchical_spatiotemporal_framework}

To manage high-velocity, time-varying pluvial rainfall fields without the computational overhead of heavy binary database containers (e.g., HDF5 or NetCDF), the analytical engine establishes a lightweight storage specification termed the Hierarchical IDRISI Format (HIF). This format emulates three-dimensional (3D) spatiotemporal data cubes using native two-dimensional (2D) GIS raster matrices as fundamental building blocks.

\subsubsection{Conceptual Abstraction and Volumetric Representation}
\label{subsub:conceptual_abstraction}

The framework discretizes a continuous spatiotemporal variable $\mathcal{F}(x, y, t)$ over an $N_x \times N_y$ spatial grid across $T$ discrete intervals. The HIF architecture structures this 3D volume as an ordered set of isolated 2D tensor slices:

\begin{equation}
    \mathcal{F} = \{ \mathbf{M}_1, \mathbf{M}_2, \dots, \mathbf{M}_t, \dots, \mathbf{M}_T \}
\end{equation}

where each matrix $\mathbf{M}_t \in \mathbb{R}^{N_x \times N_y}$ represents the spatial state at time-step $t$. 

This enables a \textsl{lazy loading} paradigm: the engine ingests only the frame $\mathbf{M}_t$ required for the active hydraulic step rather than pre-allocating memory for the entire timeline. Consequently, the memory footprint scales independently of the total duration $T$, satisfying the hardware constraints of real-time deployment.

\subsubsection{Components of the Spatiotemporal Dataset}

The framework organizes high-dimensional temporal datasets into three logical layers:

\begin{itemize}
    \item \textbf{Global Spatial Descriptor}: Metadata defining the unified coordinate reference system, boundaries, and resolution.
    \item \textbf{Temporal Attribute Manifest}: An indexed map linking discrete simulation steps to timestamps and temporal downscaling factors.
    \item \textbf{Temporal Frame Repository}: A contiguous collection of raw matrix snapshots capturing the spatial state variables per time-step.
\end{itemize}

\subsubsection{Chronological Event Synchronization and Spatial Downgrading}
\label{subsub:chronological_synchronization}

To align static physical domain attributes with dynamic boundary forcing fields (e.g., precipitation) without enforcing rigid simulation intervals, the framework uses a timestamp-driven validation loop. A dynamic environmental layer maps its state transitions to an ordered vector of timestamps $\mathcal{T}_{\text{layers}} = \{ \tau_0, \tau_1, \dots, \tau_K \}$. A frame update from $k$ to $k+1$ is triggered when the integration clock $t$ crosses the chronological threshold:

\begin{equation}
    t \ge \tau_{k+1}
\end{equation}

Additionally, when input datasets operate at a coarser resolution than the computational grid by an integer scaling factor $\gamma$, the framework applies a spatial downgrading mapping $\mathcal{P}: \mathbb{N}^2 \rightarrow \mathbb{N}^2$:

\begin{equation}
    \mathcal{P}(r, c) = \left( \lfloor r / \gamma \rfloor, \, \lfloor c / \gamma \rfloor \right)
\end{equation}

Upon a temporal frame update, the engine executes this cell-by-cell coordinate projection—utilizing an intermediate resampled buffer if $\gamma > 1$—to ensure deterministic spatial alignment without runtime memory overhead.


\subsection{The GIS-to-Tensor Pipeline: From Rasters to Multi-Channel Cubes}

The transition from environmental raw data to a simulation-ready state follows a specialized Extract, Transform, Load (ETL) framework. 
This process ensures that diverse environmental layers—specifically topography, land cover, and rainfall—are spatially and temporally 
synchronized to establish a mathematically consistent initial state for the cellular automata model.

\paragraph{A. Spatial Alignment and Domain Standardization}

Raw environmental datasets typically exhibit heterogeneous spatial resolutions, diverse geometric extents, and distinct geographic 
projections. To ensure spatial consistency across the simulation domain, the pipeline executes two primary geometric transformations:

\begin{itemize}
    \item \textbf{Coordinate Reference System (CRS) Standardization}: All geospatial layers are projected to a common metric coordinate 
    system (UTM 30N). This uniformity is mathematically required to ensure that Euclidean distance calculations within the Moore 
    neighborhood remain valid and isotropic across the spatial grid.

    \item \textbf{Resampling and Spatial Discretization}: Spatial layers are resampled to a uniform grid extent of dimensions $H \times W$. 
    The pipeline employs bilinear interpolation for continuous environmental variables (such as elevation models) and nearest-neighbor 
    interpolation for discrete thematic categories (such as land use maps) to preserve categorical boundaries.
\end{itemize}

\paragraph{B. Tensorization: The Multi-Channel Spatial Representation}

Once spatially aligned, the heterogeneous physical layers are mapped into a unified algebraic structure. The localized state of the river 
basin is represented as a 3D tensor of shape $(C, H, W)$, defined as follows:

\begin{itemize}
    \item \textbf{Channels ($C$)}: Represents the $n$ distinct physical and environmental information layers dictated by the $m:n-CA^k$ 
    model architecture (e.g., Channel 0: Water Depth, Channel 1: Land Use, Channel 2: Topography).

    \item \textbf{Height and Width ($H, W$)}: Represent the dimensions of the discretized spatial grid.
\end{itemize}

By structuring the spatial domain as a multi-channel tensor, the geographical space is treated as a single mathematical object. This 
formulation allows the cellular automata transition rules to be expressed as global tensor operations and matrix shifts, shifting the 
computational paradigm from local, independent cell evaluations to unified spatial transformations.


\subsection{Data Pipeline Design and Pre-processing Framework}

The core of the flood simulation system relies on its ability to process vast amounts of heterogeneous geospatial data accurately and 
efficiently. The data pipeline is conceived as a centralized conceptual orchestrator responsible for the ingestion, transformation, and 
provisioning of the physical parameters required by the hydraulic model.

The primary methodological objective is to decouple data acquisition from the underlying simulation logic. By isolating the data ingest 
mechanics, the system guarantees that variations in external data sources or file formats do not compromise the mathematical stability of 
the simulation engine. This framework is structured upon three fundamental pillars (Figure \ref{fig:high-level_system_context}):

\begin{enumerate}
    \item \textbf{Spatial Integrity}: Ensuring that spatial coordinate reference systems, geometric extents, and resolutions remain 
    mathematically consistent across all processed data layers.

    \item \textbf{Scalability}: Allowing the seamless inclusion of new physical variables (e.g., soil moisture, infiltration capacity) 
    without requiring structural changes to the core simulation flow.

    \item \textbf{Automation}: Minimizing manual intervention during the preprocessing stages of raw geospatial datasets to ensure workflow 
    reproducibility.
\end{enumerate}

\begin{figure}[htbp]
    \centering
    \includegraphics[width=0.85\textwidth]{chapter_3/high-level_system_context}
    \caption{High-Level System Context}
    \label{fig:high-level_system_context}
\end{figure}

\subsubsection{Data Categorization: Static vs. Dynamic Layers}

To optimize domain state updates and operational memory management, the methodology categorizes geospatial information into two distinct 
functional types based on their temporal dynamics during a simulation event:

\begin{itemize}
    \item \textbf{Static Layers}: These represent environmental characteristics that remain invariant throughout the simulated event. 
    Examples include the Digital Elevation Model (DEM) and Land Use. These layers are processed a single time during 
    initialization and serve as the immutable topographic canvas of the simulation domain.

    \item \textbf{Dynamic Layers}: These comprise time-varying boundary conditions and state variables that actively drive the transition 
    functions of the model. Examples include spatialized Rainfall Intensity. The methodology must formalize a sequential state retrieval 
    mechanism, ensuring that the simulation engine receives the exact temporal "snapshot" corresponding to each discrete time step.
\end{itemize}

This conceptual dichotomy allows the pipeline to separate the operational processing paths for each layer type, maximizing data-flow 
efficiency before execution.

\subsubsection{Data Flow Logic and Validation Principles}

The progression of data through the preprocessing pipeline is governed by a deterministic lifecycle, ensuring that the simulation engine 
receives a spatially and temporally synchronized state of the environment.

\paragraph{A. Configuration-Driven Parameterization}

The pipeline initializes via a centralized metadata control plane. This approach separates the layer definitions—which are structurally 
bound to the pipeline's logic for consistency—from their operational parameters (e.g., geographic boundaries, spatial resolutions, and 
specific interpolation models). This externalization ensures scientific reproducibility, enabling identical pipeline logic to be deployed 
across distinct geographic catchments by adjusting numerical and spatial constants.

The control plane orchestrates three main dimensions:

\begin{enumerate}
    \item \textbf{Global Simulation Metadata}: Defines temporal boundaries (start/end timestamps) and domain constraints to synchronize 
    all incoming layers within a uniform space-time window.

    \item \textbf{Layer Parameterization}: Dispatches execution parameters to individual layers, including target resolutions, spatial 
    boundaries, and interpolation models (such as variogram specifications for rainfall spatialization).

    \item \textbf{Visualization Diagnostics}: Identifies which physical layers require graphical rendering for research and diagnostics, 
    isolating them from background computational layers.
\end{enumerate}

\paragraph{B. Spatial Synchronization and Dependency Mapping}

Before executing any raster transformations, the pipeline must validate the logical structure of the data dependencies and enforce spatial 
convergence. This validation is modeled using Graph Theory principles:

\begin{itemize}
    \item \textbf{Directed Acyclic Graph (DAG) Resolution}: The operational dependencies between physical layers (e.g., a rainfall layer 
    requiring an elevation model for spatial clipping) are modeled as nodes and directed edges. The pipeline enforces a topological sort 
    to guarantee that circular dependencies are nonexistent and that each component is initialized only when its structural prerequisites 
    are met.

    \item \textbf{Master Grid Alignment}: To mitigate spatial misalignment errors to which cellular automata are highly sensitive, the 
    framework enforces a "Master Grid" paradigm. The primary topographic layer (DEM) establishes the baseline spatial constraints. All 
    secondary layers (e.g., roughness, rainfall) are validated, resampled, and clipped against this reference to guarantee exact overlap 
    in Coordinate Reference System (CRS), pixel resolution, and bounding boxes.
\end{itemize}

\paragraph{C. Sequential Execution}

The pipeline adheres to a strictly sequential execution model, processing each geospatial layer in a deterministic order derived from the 
topological sorting. Each environmental variable is isolated within its own functional generator, encapsulating the transformation rules
to preserve strict modular boundaries.
